% ---------------------------------------------------------------------------
% LakeGen extension: experimental evaluation (RQ1-RQ4) and results.
% Drop-in sections for the llncs source; include with \input{evaluation_section}
% after \input{selection_section}, to which it refers (sec:retrieval,
% sec:retrieval-comparison, sec:selection, tab:selection-tools, sec:solr).
% Needs \usepackage{tikz} for the waterfall figure. Cites pneuma and
% ji2025target (already used) and li2023bird (added to retrieval_section.bib).
% The section levels follow the other drop-ins; demote them if the paper puts
% the evaluation under one \section.
%
% GENERATED: the tables and the figure are filled from the run data by
%   PYTHONPATH=.:src .venv/bin/python experiments/paper_analysis/question_meta.py
%   PYTHONPATH=.:src .venv/bin/python experiments/paper_analysis/metrics.py
%   PYTHONPATH=.:src .venv/bin/python experiments/paper_analysis/failures.py
%   PYTHONPATH=.:src .venv/bin/python experiments/paper_analysis/latex.py
%   PYTHONPATH=.:src .venv/bin/python experiments/paper_analysis/rerank_compare.py
%   PYTHONPATH=.:src .venv/bin/python experiments/paper_analysis/section_tables.py
%   python3 experiments/paper_analysis/build_section.py
% The numbers quoted in the text are in results/section_numbers.json. Edit
% experiments/paper_analysis/evaluation_section.template.tex, not the generated
% paper/evaluation_section.tex.
% ---------------------------------------------------------------------------

\section{Experimental Evaluation}
\label{sec:evaluation}

We evaluate the retrieval strategies of Section~\ref{sec:retrieval}, the two
selectors of Section~\ref{sec:selection} and the code they lead to. For every
question we follow the answer through retrieval, selection and coding, and we
attribute each failure to the first stage that loses it.

\subsection{Experimental Questions}
\label{sec:rq}

\begin{description}
\item[RQ1 (retrieval).] How completely, and how high, does each retrieval
strategy place the tables a question needs among its candidates, and does a
second-stage re-ranker improve their order?
\item[RQ2 (selection).] Given the same candidates, how accurately do the LLM
selector and Jev identify the subset of tables that the question requires?
\item[RQ3 (portals and languages).] How do retrieval and selection vary across
portals, platforms and catalogue languages, and do the findings hold on two
external benchmarks?
\item[RQ4 (propagation).] How do retrieval and selection errors propagate to
the final answer, and how does this depend on the number of tables, on how
they combine, and on the difficulty of the question?
\end{description}

% NOTE for the authors, not for print: align the wording of RQ1-RQ4 with the
% introduction.

\subsection{Datasets}
\label{sec:datasets}

Table~\ref{tab:datasets} describes the lakes. On each of the five portals we
use 100 questions generated with the OrQA workflow and sampled with a fixed
seed: a third each of easy, medium and hard questions, 66 over one table and
34 over several. OrQA labels each multi-table question by how its tables
combine: a \emph{join} on a key, a \emph{union} of tables with the same schema
(for instance the editions of one series), or a \emph{join-correlation}, a
join whose question relates columns of the two tables. Each question comes
with a reference program and its result. Over the 500 portal questions, the
reference programs aggregate in 86\% of the cases, filter in 70\%, derive new
columns in 50\%, sort in 48\%, group in 23\% and clean values in 8\%.

Two external benchmarks test whether the findings carry over to data that was
not generated by our workflow. TARGET~\cite{ji2025target} pools the 75 tables
of the eleven development databases of BIRD~\cite{li2023bird} into one lake.
We draw 100 of its questions in proportion to BIRD's difficulty levels; three
in four need a join. The Chicago lake is the Chicago Open Data benchmark of
Pneuma~\cite{pneuma}, 801 tables with questions about their content. Its gold
standard lists, for each question, every table from which the answer can be
computed (up to 80), and any one of them suffices.

\begin{table}[t]
\caption{Lakes and questions. S/M: questions over a single table or
multiple tables (Chicago: any one of up to 80 gold alternatives). J/U/JC: multi-table questions combining their tables by a join, a
union or a join-correlation. E/M/H: easy, medium and hard questions (BIRD's
simple, moderate and challenging on TARGET). Languages: es Spanish, ca
Valencian, fr French, it Italian, en English.}
\label{tab:datasets}
\centering
\footnotesize
\setlength{\tabcolsep}{3pt}
\begin{tabular}{llllrccc}
\hline
Lake & Platform & Lang. & Questions & Tables & S/M & J/U/JC & E/M/H \\
\hline
Valencia & CKAN & es, ca & OrQA & 122 & 66/34 & 9/16/9 & 34/33/33 \\
Paris & OpenDataSoft & fr & OrQA & 471 & 66/34 & 13/14/7 & 34/33/33 \\
Bologna & OpenDataSoft & it & OrQA & 677 & 66/34 & 16/16/2 & 34/33/33 \\
UK & CKAN & en & OrQA & 15{,}205 & 66/34 & 14/18/2 & 34/33/33 \\
NYC & Socrata & en & OrQA & 2{,}673 & 66/34 & 16/14/4 & 34/33/33 \\
\hline
TARGET & BIRD & en & BIRD & 75 & 25/75 & 75/0/0 & 61/30/9 \\
Chicago & Socrata & en & Pneuma & 801 & any & -- & -- \\
\hline
\end{tabular}
\end{table}

\subsection{Compared Retrieval Methods}
\label{sec:methods}

We compare the six strategies of Table~\ref{tab:methods}, described in
Section~\ref{sec:retrieval}. They differ in what they read and in what an LLM
contributes at query time. Hybrid content search is not part of this
evaluation.

\begin{table}[t]
\caption{Compared retrieval methods: the evidence each one reads (catalogue
metadata, embeddings, cell values; Pneuma reads a sample of rows), whether it
needs an index built outside Solr, and the role of an LLM.}
\label{tab:methods}
\centering
\footnotesize
\setlength{\tabcolsep}{3pt}
\begin{tabular}{lcccc p{3.0cm}}
\hline
Method & Metadata & Embed. & Cells & Ext.\ index & LLM at query time \\
\hline
Keyword & $\bullet$ & & & & keywords \\
Semantic & $\bullet$ & $\bullet$ & & & -- \\
Hybrid & $\bullet$ & $\bullet$ & & & keywords \\
Content & $\bullet$ & & $\bullet$ & & keywords \\
Pneuma & $\bullet$ & $\bullet$ & sample & $\bullet$ & relevance judge (narrations offline) \\
Pneuma-Seeker & $\bullet$ & $\bullet$ & $\bullet$ & $\bullet$ & judge, entity extraction (narrations offline) \\
\hline
\end{tabular}
\end{table}

\subsection{Second-Stage Re-ranking}
\label{sec:rerank-setup}

The re-ranker reorders the 20 candidates of a search and never adds or
removes one. To isolate its effect we evaluate every strategy in two
settings on the same candidates: \emph{no re-ranking}, the order of the
first-stage score, and \emph{Cohere re-ranking}, the order given by Cohere
Rerank on the text described in Section~\ref{sec:retrieval}. In production
keyword, semantic and hybrid search are re-ranked and the other three are not.
For each strategy the setting that production does not use is computed on the
same saved candidates. Since the candidate set is the same in both settings,
re-ranking cannot change R@20 or retrieval success, only the order.

\subsection{Metrics}
\label{sec:metrics}

\paragraph{Retrieval success and selection exactness.}
We keep two questions apart. \emph{Retrieval success} asks whether every
table the question requires entered the candidate set. \emph{Selection
exactness} asks whether the selector then picked exactly those tables. If the
gold tables are $\{A,B\}$ and retrieval returns $\{A,C,D,E,F\}$, then R@5 is
0.5 and the candidate set is incomplete: no selector can produce the answer. If
retrieval returns $\{A,B,C,D,E\}$, then R@5 is 1 and P@5 is 0.4, and the
selection may still be exact if the selector keeps $\{A,B\}$.

\paragraph{Retrieval.}
R@$K$ is the fraction of the gold tables among the first $K$ candidates
($K=1,5,10,20$). On a multi-table question R@1 cannot exceed the inverse of
the number of gold tables. \emph{Cpl@20}, retrieval success, is 1 when every
gold table is among the 20 candidates. We also report the mean reciprocal rank
of the first gold table (MRR) and nDCG@10 with binary relevance. P@$K$ is
omitted, since with one or two gold tables per question it mostly reflects
$K$.

\paragraph{Selection.}
For a selected set $S$ and gold set $G$: set recall $|S\cap G|/|G|$, set
precision $|S\cap G|/|S|$ (0 for an empty selection), their F1, exact set
match ($S=G$), \emph{answerable} ($G \subseteq S$: the selection contains every
table the reference program reads), the unnecessary-table rate (the share of
selections with a table outside $G$), and the empty-selection rate.

\paragraph{Downstream.}
For the code run on a selection: \emph{execution success}, the code ran and
returned a result; \emph{exact match}, the result equals the reference result,
up to representation and a relative tolerance of $10^{-6}$ on numbers;
\emph{lenient match}, the same rows when columns are renamed or a proportion is
given as a percentage; and \emph{answer correctness}, an exact match or a
result the semantic code judge accepts as an alternative correct answer.
\emph{Final QA success} is answer correctness over all questions, with an
empty selection counted as a failure.

On Chicago, where any gold alternative suffices, retrieval success, R@$K$ and
an answerable or exact selection all mean that at least one alternative is
present.

\subsection{Protocols}
\label{sec:protocols}

On the five portals we use the \emph{same-context} protocol of
Section~\ref{sec:selection}. Retrieval runs once per question and strategy,
with the keywords given in the benchmark. Both selectors choose from the same
20 candidates: the LLM selector in a single call without tools, Jev with every
candidate. Each selected set of tables is then coded with the requirement plan
switched off, and the outcome is attributed to every strategy and selector
that chose that set. A set coded more than once counts with its mean outcome.
Every set chosen by the LLM selector was coded, as were 81\% of those chosen
by Jev; Jev's uncoded sets are left out of its downstream rates. On TARGET and
Chicago we run the whole pipeline in the divided architecture, once per
strategy and selector. The Retriever agent writes the keywords, the LLM
selector is the agent with the tools of Table~\ref{tab:selection-tools}, and
Jev receives the top 8 candidates. We report the final retrieval attempt, the
final selection and the outcome of the run. Every value comes from a single
run.

\subsection{Configuration}
\label{sec:configuration}

Tables~\ref{tab:setup} and~\ref{tab:retrieval-parameters} give the models,
software and parameters. Every LLM call uses GPT-OSS 120B at the same
temperature, through two providers, Oracle Cloud Infrastructure (OCI)
Generative AI and Together AI.

\begin{table}[t]
\caption{Models, software and hardware.}
\label{tab:setup}
\centering
\small
\begin{tabular}{lp{7.6cm}}
\hline
Component & Setting \\
\hline
Keyword generation & GPT-OSS 120B (Retriever agent). The same-context protocol uses the benchmark's keywords, written by the OrQA generator (GPT-OSS 120B) \\
Selection (LLM) & GPT-OSS 120B \\
Integration and coding & GPT-OSS 120B, at most 3 revision rounds \\
Code judge & GPT-OSS 120B \\
Benchmark generation & OrQA with GPT-OSS 120B on OCI (portals); BIRD and Pneuma questions as released (TARGET, Chicago) \\
LLM providers & OCI Generative AI and Together AI; at most 16{,}000 output tokens, 128K context \\
Temperature & 0.1 for every GPT-OSS 120B call \\
Random seed & 0 for \textsc{LakeGen}'s own random generators; 42 for sampling the questions. The LLM APIs are not seeded \\
Embedding model & Cohere Embed v4.0 (\texttt{cohere.embed-v4.0}), 1{,}024 dimensions, document and query input types \\
Re-ranker & Cohere Rerank v4.0 fast (\texttt{cohere.rerank-v4.0-fast}) \\
Pneuma & version 0.0.4 with the corrections of Section~\ref{sec:retrieval}; narrations and judge by GPT-OSS 20B on OCI (narrations at temperature 0.7, judge at 0 with seed 42); Cohere Embed v4.0 \\
Jev & \texttt{jev-1.13.0} (TypeSafe API) \\
Solr & 9.10.1 (Lucene 9.12.3, OpenJDK 17) \\
DuckDB & 1.5.2 (1.5.5 in the Pneuma service) \\
Python & 3.13.13 (3.11.2 in the Pneuma service) \\
Hardware & Intel Core i9-13900KF (24 cores, 32 threads), 128\,GB RAM, NVMe SSD, Debian 12. No model runs locally \\
\hline
\end{tabular}
\end{table}

\begin{table}[t]
\caption{Retrieval and selection parameters. $k$: candidates returned by a
search.}
\label{tab:retrieval-parameters}
\centering
\small
\begin{tabular}{lp{8.4cm}}
\hline
Stage & Parameters \\
\hline
All strategies & $k=20$; the LLM agent sees the first 10 and Jev the first 8 (pipeline), or both see all 20 (same-context) \\
Keyword & Solr extended DisMax, AND operator, BM25; field weights: title 4, tags and column names 3, title $n$-grams, description and column labels 2, others 1; cascade of up to 5 keyword sets; up to 2 new searches after an empty result \\
Semantic & KNN on an HNSW index, cosine similarity \\
Hybrid & weighted fusion, $\alpha=0.25$ (lexical weight), $c=5$ ($c \cdot k = 100$ candidates per branch), min-max normalisation, score 0 for a missing branch \\
Content & first pass over all tables; second pass on the best 250 (first 40 text columns, 100{,}000 rows each); 25 tables probed on 1{,}000 rows; terms outside $K$ weighted 0.25; at most 3 keywords covered per table; 3 sample rows \\
Pneuma & $k=20$, $n=5$, $\alpha=0.25$; timeout 600\,s \\
Pneuma-Seeker & content weight 0.5, merge by filling Pneuma's free slots (reference implementation), no enumeration; evidence weights title 3, column name 2, cell 1 \\
Re-ranking & depth 20; in production on keyword, semantic and hybrid only \\
LLM selector (agent) & at most 5 inspected candidates, 8 after one expansion of 5; 16 iterations, 12 tool calls and 300\,s per attempt; at most 3 attempts \\
Jev & threshold $\tau=0.5$; back to retrieval below 0.15; requests capped at 80\,KB; profiles on the first 200{,}000 rows \\
\hline
\end{tabular}
\end{table}

% NOTE for the authors, not for print. (1) The portal results (same-context
% protocol: benchmark keywords, tool-free LLM selector) and the TARGET/Chicago
% results (whole pipeline: agent keywords, agent selector) are not directly
% comparable, which is why they are reported apart. (2) UK was measured on the
% Solr core with the 3,392 "CSV" titles, before the fix of sec:solr; the fixed
% core has no vectors yet. (3) Providers: of the portal codings 3,887 ran on OCI
% and 1,074 on Together; of the 2,400 pipeline runs about 590 ran on OCI before
% the switch to Together, then 894 on Together and 379 on OCI. (4) Check that
% the OrQA generator never fell back to its configured backup model (Groq
% Llama 3.3 70B).

\section{Results}
\label{sec:results}

\subsection{Retrieval Performance (RQ1)}
\label{sec:results-retrieval}

\begin{table}[p]
\caption{Retrieval with each strategy's production setting (re-ranking on
keyword, semantic and hybrid search), 100 questions per lake. Above the double
line: the five portals (same-context protocol, benchmark keywords) and their
mean. Below: $\ddagger$ TARGET and Chicago (whole-pipeline protocol, keywords
written by the agent, mean of the runs with both selectors; on Chicago one
gold alternative counts). Cpl@20: every gold table among the 20 candidates.
Best value per lake and column in bold.}
\label{tab:rq1}
\centering
\small
\setlength{\tabcolsep}{3pt}
\begin{tabular}{ll|cccc|cc|c}
\hline
Portal & Strategy & R@1 & R@5 & R@10 & R@20 & MRR & nDCG@10 & Cpl@20 \\
\hline
%%RQ1_ROWS%%
\end{tabular}
\end{table}

Table~\ref{tab:rq1} reports retrieval on the portals. Averaged over them, the
hybrid of BM25 and KNN had the highest value on every measure: all gold tables
were among its 20 candidates for 96\% of the questions, and its MRR was 0.85.
It led on four portals. On UK the two Pneuma rankings retrieved more questions
completely (0.94 against 0.86). Pneuma and Pneuma-Seeker came close to the
hybrid on retrieval success (Cpl@20 0.93) but ranked the gold tables lower
(MRR 0.71), whereas semantic search ranked them higher (MRR 0.80) but
retrieved fewer questions completely (0.89). Content search reached 0.71, but
0.93 on NYC. Keyword search reached 0.62, and its
recall did not grow beyond the fifth candidate (0.67 at 5, 10 and 20): its AND
queries return short lists, so a gold table missing from them is not retrieved
at all. The two benchmarks, run through the whole pipeline, keep the picture:
semantic, hybrid and Pneuma-Seeker retrieval led on TARGET (Cpl@20
0.96--0.97), Pneuma on Chicago (0.96), and keyword search was last on both.

\begin{table}[t]
\caption{Retrieval by question scope and by how the tables combine, pooled
over the five portals: R@10 and Cpl@20.}
\label{tab:scope}
\centering
\small
\setlength{\tabcolsep}{3.5pt}
\begin{tabular}{l|cc|cc|cc|cc|cc}
\hline
 & \multicolumn{2}{c|}{Single} & \multicolumn{2}{c|}{Multi} & \multicolumn{2}{c|}{Join} & \multicolumn{2}{c|}{Union} & \multicolumn{2}{c}{Join-corr.} \\
Strategy & R@10 & Cpl & R@10 & Cpl & R@10 & Cpl & R@10 & Cpl & R@10 & Cpl \\
\hline
 & %%SCOPE_N%% \\
\hline
%%SCOPE_ROWS%%
\hline
\end{tabular}
\end{table}

Separating the questions by scope (Table~\ref{tab:scope}) shows where the
strategies differ. On single-table questions keyword search retrieved the gold
table nearly as often as semantic search (Cpl@20 0.89 against 0.92). On multi-table questions it
retrieved every gold table for 9\% of them, because one conjunctive query
rarely matches several tables. The hybrid kept a Cpl@20 of 0.92 on
multi-table questions. Unions were the hardest case for complete retrieval
(0.86 for the hybrid, 0.79 for semantic search), since all the partitions of
a series must enter the candidates, and they compete with one another for the
same 20 slots.

\subsection{Effect of Second-Stage Re-ranking (RQ1)}
\label{sec:results-rerank}

\begin{table}[t]
\caption{Re-ranking off and on, on the same 20 candidates, pooled over the five
portals. $\dagger$: the production setting. Better value of each pair in bold.
Last column: portals on which re-ranking raised the MRR.}
\label{tab:rerank}
\centering
\small
\begin{tabular}{ll|ccc|cc|c}
\hline
Strategy & Re-ranking & R@1 & R@5 & R@10 & MRR & nDCG@10 & Portals \\
\hline
%%RERANK_ROWS%%
\end{tabular}
\end{table}

Table~\ref{tab:rerank} isolates the re-ranker. On these portals it improved
the order of every strategy, and most where the first-stage order was weakest.
For content search the MRR rose from 0.55 to 0.72 and R@1 from 0.35 to 0.51.
For Pneuma and Pneuma-Seeker, whose service returns an order without scores,
the MRR rose from 0.71 to 0.83. Keyword search gained 0.07 MRR. The hybrid,
already the best ordered, gained less than 0.01, and on NYC the re-ranker
lowered the MRR of semantic and hybrid search (by 0.005 and 0.03). These gains concern
the order only. In the same-context protocol both selectors saw all 20
candidates, so the order could matter only through the selectors' attention.
In the pipeline the agent sees the first 10 and Jev the first 8, so a better
order can bring a gold table into view. The re-ranker reads catalogue
metadata, so it may help where the metadata distinguishes the candidates and
override content evidence where it does not. On these portals the first
effect dominated, although production keeps content search and the Pneuma
rankings unre-ranked. A single run per setting does not settle whether the
smaller differences are significant.

\subsection{Table Selection (RQ2)}
\label{sec:results-selection}

\begin{table}[t]
\caption{Selection on the five portals (same-context protocol), pooled per
strategy. Rec./Prec.: set recall and precision. Exact\,|\,Ret.: exact
selection among questions with every gold table in the candidates. Ans.:
selection containing every gold table. Unn.: selections with a table outside
the gold set. Empty: empty selections.}
\label{tab:selection}
\centering
\small
\setlength{\tabcolsep}{3pt}
\begin{tabular}{ll|ccc|cc|ccc}
\hline
Strategy & Selector & Rec. & Prec. & F1 & Exact & Exact\,|\,Ret. & Ans. & Unn. & Empty \\
\hline
%%SEL_ROWS%%
\end{tabular}
\end{table}

Table~\ref{tab:selection} compares the selectors on identical candidates. The
LLM selector matched the gold set exactly more often (0.62 against 0.58 over
all strategies, 0.70 against 0.65 with the hybrid), while Jev recovered more
of the gold tables (set recall 0.77 against 0.71) and more often selected all
of them (answerable 0.71 against 0.65). The two trade errors: Jev included a
table outside the gold set in 31\% of its selections against 24\% for the
LLM, and the LLM returned an empty selection more often (9\% against 6\%). Even
when retrieval had succeeded, an exact selection was not guaranteed: with
the hybrid it reached 0.73 for the LLM selector and 0.68 for Jev. Selection
therefore loses more questions than retrieval with the strongest strategies,
which motivates measuring the two stages separately.

\subsection{Cross-Portal and Multilingual Analysis (RQ3)}
\label{sec:results-cross}

\begin{table}[t]
\caption{Per lake, with the hybrid as retrieval strategy: Cpl@20 of keyword,
semantic and hybrid search, the strategy with the highest Cpl@20, exact
selection and final QA success per selector. $\ddagger$: whole-pipeline
protocol; on Chicago, ``exact'' means that a gold alternative was
selected.}
\label{tab:cross}
\centering
\scriptsize
\setlength{\tabcolsep}{2.5pt}
\begin{tabular}{lll|ccc|l|cc|cc}
\hline
 & & & \multicolumn{3}{c|}{Cpl@20} & & \multicolumn{2}{c|}{Exact sel.} & \multicolumn{2}{c}{Final QA} \\
Lake & Platform & Lang. & Kw. & Sem. & Hyb. & Best Cpl@20 & LLM & Jev & LLM & Jev \\
\hline
%%CROSS_ROWS%%
\hline
\end{tabular}
\end{table}

Table~\ref{tab:cross} summarises each lake. The catalogue language did not
order the portals: semantic search, which uses a multilingual model, retrieved
every gold table for 0.94 of the Valencia and 1.00 of the Paris questions,
against 0.94 on NYC and 0.76 on UK. What set UK apart was the state of its
catalogue (titles reduced to ``CSV'', Section~\ref{sec:solr}) and its size
(15{,}205 tables), not its language, and UK had the lowest selection rates for
both selectors. Keyword search, whose stemmer is the English one on every
core, did not fare worse on the three non-English catalogues (Cpl@20
0.61--0.67) than on NYC (0.72). Multilingual mismatches, a keyword in another
language than the catalogue, caused 21 of the 881 retrieval failures, 19 of
them on Valencia, whose catalogue mixes Spanish and Valencian.

On the external benchmarks the retrieval ranking held: semantic and hybrid
search retrieved every gold table for 96\% of the TARGET questions, and the
Pneuma rankings did best on Chicago, the benchmark on which Pneuma was
developed (Cpl@20 0.96). Selection did not transfer the same way. On TARGET,
where three questions in four need a join, the LLM agent selected the exact
set for 15\% of the questions against 66\% for Jev, and Jev doubled the final
QA success (0.39 against 0.19). On Chicago the two selectors were close.

\subsection{Propagation to Table Question Answering (RQ4)}
\label{sec:results-propagation}

\begin{figure}[t]
\centering
%%WATERFALL_FIGURE%%
\caption{Retrieval-to-answer waterfall on the five portals (500 questions),
with hybrid retrieval and the LLM selector: questions whose candidates contain
every gold table (retrieved), whose selection contains every gold table
(answerable), whose code then ran (executed), and whose answer was correct.
Each bar keeps only the questions that passed every earlier stage.}
\label{fig:waterfall}
\end{figure}

\begin{table}[t]
\caption{Retrieval-to-answer waterfall per lake with hybrid retrieval: share of
questions (\%) that passed each stage and every earlier one. Exact sel.: exact
selection, for reference. Off path: correct answers although a stage failed.}
\label{tab:waterfall}
\centering
\small
\begin{tabular}{l|cccc|cc}
\hline
Lake (selector) & Retrieved & Answerable & Executed & Correct & Exact sel. & Off path \\
\hline
%%WATERFALL_ROWS%%
\hline
\end{tabular}
\end{table}

Figure~\ref{fig:waterfall} and Table~\ref{tab:waterfall} follow the questions
from retrieval to the answer. With the hybrid, retrieval lost 4\% of the portal
questions, selection a further 23\%, execution 4\%, and the code 45\%: 23\%
of the questions ended with a correct answer. Pooled over all six strategies,
retrieval lost 16\% of the questions, selection 19\%, execution 3\%, and the
code about 40\%.
Errors propagate almost without recovery. When a gold table was missing from
the candidates, 2\% of the portal questions were still answered correctly,
against 30\% when retrieval succeeded. When the selection contained every gold
table, 33\% were answered correctly, against 15\% for the other questions whose
retrieval had succeeded. TARGET and Chicago show the same pattern (2\% and 4\%
after a retrieval failure), and on TARGET the waterfall makes the selector's
role visible: with the same retrieval, 42\% of the questions reached an
answerable selection with the LLM agent and 79\% with Jev.

\begin{table}[t]
\caption{Hybrid retrieval on the five portals by difficulty, scope and
combination of tables: retrieval, selection F1, exact selection and final QA
success per selector (L: LLM, J: Jev).}
\label{tab:difficulty}
\centering
\small
\setlength{\tabcolsep}{3.5pt}
\begin{tabular}{lr|cc|cc|cc|cc}
\hline
 & & & & \multicolumn{2}{c|}{Sel.\ F1} & \multicolumn{2}{c|}{Exact sel.} & \multicolumn{2}{c}{Final QA} \\
Questions & $n$ & R@10 & Cpl@20 & L & J & L & J & L & J \\
\hline
%%DIFF_ROWS%%
\hline
\end{tabular}
\end{table}

Table~\ref{tab:difficulty} breaks the portal results down. Difficulty, as
OrQA assigns it, barely affected retrieval and selection: R@10 was 0.95--0.97
and exact selection by the LLM selector 0.64--0.73 at every level. It
affected the answer: final QA success fell from 0.36 on easy questions to 0.23
on medium and 0.21 on hard ones, which suggests that OrQA's difficulty
measures the computation more than the search for tables. The number of tables
mattered at every stage. Multi-table questions had a lower Cpl@20 (0.92
against 0.98), a lower exact selection (0.56 against 0.77 for the LLM
selector) and half the final QA success (0.16 against 0.33). Among them,
unions were the hardest to retrieve completely, joins and join-correlations
the hardest to select, and join-correlations, with 24 questions, the hardest
to answer (0.08 and 0.12). On TARGET the pattern differs with BIRD's
difficulty levels and with the selector: Jev answered 38\% of the simple and
37\% of the moderate questions, the LLM agent 25\% and 7\%.

\begin{table}[t]
\caption{Downstream measures, pooled over the six strategies. Coded: runs with
a coded non-empty selection. Executed, Exact, Lenient, Correct: over those
runs. Correct\,|\,Ans.: over runs whose selection contained every gold table.
Final QA: over all runs. Exec.: executed.}
\label{tab:downstream}
\centering
\footnotesize
\setlength{\tabcolsep}{3pt}
\begin{tabular}{lr|cccc|c|c}
\hline
Setting & Coded & Exec. & Exact & Lenient & Correct & Corr.\,|\,Ans. & Final QA \\
\hline
%%DOWN_ROWS%%
\hline
\end{tabular}
\end{table}

Table~\ref{tab:downstream} separates the downstream measures. On the portals
the code ran in 88--91\% of the runs but matched the reference in 27--29\%,
and only a third of the runs with every gold table selected were answered
correctly. Leniency mattered on the portals: 12--13 more points of the runs
returned the reference rows with renamed columns or a proportion given as a
percentage, while on TARGET and Chicago it added at most 3 points. The
semantic judge did not add correct answers beyond the exact matches: in these
runs, every answer it accepted also matched the reference. On TARGET the LLM
agent's code ran in only 59\% of the runs, against 81\% after Jev's
selections, since a selection that lacks a table needed for a join leaves the
coder without the columns it requires.

\subsection{Failure Analysis}
\label{sec:failure-analysis}

We attribute each failed run to the first stage that lost the answer: a
\emph{retrieval failure} leaves a gold table outside the candidates, a
\emph{selection failure} has every gold table among the candidates but a
selection that is not exact (on Chicago, contains no gold alternative), and a
\emph{downstream failure} has an exact selection but a wrong answer. Each
failure receives the first category, in the order of Table~\ref{tab:failures},
whose rule holds, and the facts behind every assignment are stored with it. The
rules are the following.

\emph{Retrieval} (on the first missing gold table). Missing metadata: a
generic title such as ``CSV'', a description under 20 characters, or no column
names on a portal that publishes them. Wrong year or period: a retrieved
non-gold table has the same title once years, months and quarters are
removed. Wrong geographic or entity scope: a retrieved non-gold table's title
differs from the gold title in at most four words, shares at least half of
them, and the difference is a geographic term or a name used in the question.
Multilingual mismatch: a keyword that matches nothing in the lake's catalogue
but matches the English NYC catalogue. Entity or value only in cells: a
keyword or a value named in the question occurs in the gold table's text
cells (first 200{,}000 rows) but not in its catalogue entry. Overly restrictive
AND query: for keyword and hybrid search, the gold entry matches some keywords
but not all. Terminology mismatch: it matches none of them. The remaining
failures matched the keywords but were ranked below the cut-off.

\emph{Selection.} Relevant candidate not inspected: a gold table ranked beyond
what the selector examines (10 candidates for the agent, 15 after an
expansion, 8 for Jev; pipeline runs only). Join relationship not recognised:
some but not all gold tables of a multi-table question were selected.
Unnecessary table included: every gold table was selected, together with
another one. Column-level ambiguity: a selected non-gold table shares at least
half of its columns with a gold table that was not selected. Incorrect table
rejection: no gold table was selected, including empty selections.

\emph{Downstream.} The semantic code judge splits each question into typed
requirements (filter, aggregation, grouping, output, \dots) and marks each one
verified, failed or unknown. Execution error: no attempt ran. Incorrect
cleaning: a failed requirement whose evidence concerns casting, parsing,
missing values, duplicates or whitespace. Wrong filter: a failed filter or
scope requirement. Wrong aggregation: a failed aggregation, computation or
grouping requirement. Correct tables but wrong query: any other wrong answer.
Two cases fall outside the taxonomy: the coder rejected the selected tables,
or the judge marked no requirement as failed but could not verify the answer.

\begin{table}[t]
\caption{Failures by stage and category on the seven lakes, pooled over
strategies and selectors: count and, in parentheses, the percentage of the
lake's failures at that stage. Retrieval failures are counted once per
question and strategy on the portals, whose selectors share the candidates;
downstream failures once per coded table set on the portals and once per run
on TARGET and Chicago. Category names shortened from the text. $^{*}$: outside
the taxonomy.}
\label{tab:failures}
\centering
\scriptsize
\setlength{\tabcolsep}{2pt}
\begin{tabular}{lrrrrrrrr}
\hline
Category & VLC & PAR & BLQ & UK & NYC & TGT & CHI & Total \\
\hline
%%FAILURE_ROWS%%
\end{tabular}
\end{table}

\begin{table}[t]
\caption{One example per failure category (lake; selector where it matters).}
\label{tab:failure-examples}
\centering
\scriptsize
\begin{tabular}{lp{8.0cm}}
\hline
Category & Example \\
\hline
\multicolumn{2}{l}{\emph{Retrieval}} \\
Missing metadata & UK: ``Which strategic clinical network covered the most local authority districts on 31 December 2018?'' The gold table's catalogue title is ``CSV''. \\
Wrong year or period & Chicago: ``Which address corresponds to the condition that September is 7626?'' The 2019 edition of \emph{Libraries -- Circulation by Location} was retrieved, not the 2023 one. \\
Wrong geo./entity scope & Bologna: a question on traffic gates 39, 43, 62 and 66. Gate 51 (\emph{Capo di Lucca}) was retrieved, gate 43 (\emph{del Lavoro}) was not. \\
Multilingual mismatch & Valencia: ``\textquestiondown Cu\'antas estaciones commuter hay en Valencia?'' The keyword \emph{commuter} does not occur in a catalogue that calls them \emph{estaciones de cercan\'{\i}as}. \\
Value only in cells & TARGET: ``What is Eddie Irvine's average points scored in year 2000?'' The name occurs only in the cells of \texttt{drivers}. \\
Restrictive AND query & UK: payments ``over \pounds 25,000 made by the National Gallery in June 2022 and July 2023''. The July 2023 file matches \emph{over}, \emph{2023}, \emph{july}, \emph{payments}, not \emph{2022}, \emph{june}. \\
Terminology mismatch & TARGET: ``For the client who applied 98832 USD loan in 1996/1/3, when was his/her birthday?'' Keywords \emph{borrower}, \emph{birthday}; the table is \texttt{client}. \\
\multicolumn{2}{l}{\emph{Selection}} \\
Candidate not inspected & TARGET (LLM agent): ``State what kind of expenses that Sacha Harrison incurred?'' \texttt{member} was ranked 11th, beyond the 10 candidates shown. \\
Join not recognised & NYC (LLM): golf courses listed in both the DPR directory and the NYC Parks active list. Only one of the two lists was selected. \\
Unnecessary table & UK (Jev): average salaries of the top 10 senior civil servant posts of the Department for Transport in 2022. A second table was selected with the gold one. \\
Column-level ambiguity & UK (LLM): average bush contribution score of the cells scoring 70 or more. A table with the same columns was selected instead of the gold one. \\
Incorrect rejection & UK (LLM): the year with the most new listed buildings in Rossendale. The gold table was ranked first; nothing was selected. \\
\multicolumn{2}{l}{\emph{Downstream}} \\
Execution error & Chicago: docks available and percentage full of the stations in service. Every attempt failed on a missing column (\texttt{station name}). \\
Incorrect cleaning & TARGET: badges obtained ``on 7/19/2010 7:39:08 PM''. The code parsed the date into a timestamp and removed duplicate names, returning 3 of the 13 rows. \\
Wrong filter & TARGET: black Artifact cards with a foreign-language translation. Artifacts were filtered on \texttt{type} only, black on \texttt{colors} only. \\
Wrong aggregation & TARGET: users aged 40 among comments with score 0. A wrong join counted 31 users instead of 1{,}975. \\
Right tables, wrong query & UK: the five suppliers with the highest payments over \pounds 25,000 in two months. The ranking of the suppliers differs from the reference. \\
\hline
\end{tabular}
\end{table}

Tables~\ref{tab:failures} and~\ref{tab:failure-examples} report the categories
and an example of each. \emph{Retrieval.} Nearly half of the 881 retrieval
failures (48\%) concern questions that identify their table through a value
only its cells contain: 70\% on Chicago, whose questions are generated from
cell values, and 43--61\% on Valencia, Paris, Bologna, NYC and TARGET. The
other categories follow the habits of each catalogue. On UK, 45\% of the
retrieval failures are gold tables titled ``CSV'' and 14\% are editions of a
series for another period; on Bologna, 21\% are tables built from one template
for another electoral list or traffic gate; on TARGET, 31\% are terminology
mismatches, since its tables carry the names of database relations rather
than the words of the questions. Overly restrictive AND queries account for
11\% and multilingual mismatches for 2\%.

\emph{Selection.} Of the 2{,}177 selection failures, 32\% miss a table the
question needs to join or union and 31\% add one it does not need, and the two
selectors fail in opposite ways: missed joins are 40\% of the LLM selector's
failures and 24\% of Jev's, unnecessary tables 19\% and 44\%. On TARGET, 236
of the 293 missed joins come from the LLM agent. Column-level ambiguity, a
near-identical table chosen instead of the gold one, accounts for 41--55\% of
the failures on UK, NYC and Chicago, whose catalogues publish many tables with
the same schema. Gold tables ranked beyond the selector's window explain 4\%,
all in pipeline runs.

\emph{Downstream.} Of the 992 downstream failures, 26\% have a wrong filter,
the most frequent specific error on every lake, and 22\% another error in the
query, the largest category on Chicago (31\%). Wrong aggregations (5\%),
execution errors (3\%) and cleaning errors (2\%) are rare. In 36\% of the
cases the judge could not verify the answer, most often on TARGET and Chicago
(44\% and 38\%), and the coder rejected a correct selection in 48 cases.

Each failure receives a single category although causes may combine, the
categories are recognised by rules over recorded facts rather than read by a
person, and the counts pool strategies and selectors that fail on the same
questions. The analysis locates where answers are lost, not how often each
cause alone would prevent an answer.

% NOTE for the authors, not for print. The categories come from the rules in
% experiments/paper_analysis/failures.py; results/failures.jsonl stores, for
% each failure, the facts behind its label. Spot-check a sample before
% submission. Known weak points: "wrong geographic or entity scope" cannot tell
% a place from another named entity (on Bologna it groups traffic gates with
% electoral lists), and the cleaning rule reads the judge's evidence with a
% keyword pattern. The share of "judge could not verify" on TARGET and Chicago
% deserves a look at the judge's inputs on those lakes. The comparison table of
% sec:selection (tab:selection-results) overlaps with tab:selection and
% tab:cross; keep one of them.
